\section{Doubled-spin conventions and the finite two-channel class}
\subsection{Doubled spins}
All representation labels are nonnegative integers representing twice the ordinary spin.  Thus a doubled label $a\in\N_0$ corresponds to spin $a/2$.  For two doubled labels $a,b$, the admissible intermediate labels are
\begin{equation}
\mathcal C(a,b)=\{|a-b|,|a-b|+2,\ldots,a+b\}.
\label{eq:pairchannels}
\end{equation}

\begin{definition}[Ordered carrier]
An ordered carrier is a four-tuple
\[
J=(a,b,c,d)\in\N^4.
\]
Its two recoupling-channel sets are
\begin{equation}
E(J)=\mathcal C(a,b)\cap\mathcal C(c,d),
\qquad
G(J)=\mathcal C(b,c)\cap\mathcal C(a,d).
\label{eq:channelsets}
\end{equation}
\end{definition}

\begin{definition}[Declared two-channel class]
A carrier $J$ belongs to the declared class $\mathfrak C_6$ when:
\begin{enumerate}[label=(\alph*)]
\item $1\le a,b,c,d\le6$;
\item $a+b+c+d$ is even;
\item $|E(J)|=|G(J)|=2$;
\item every factorial argument in the finite Racah expression is nonnegative;
\item the maximum numerator factorial index appearing in the finite sum is at most $12$;
\item the anchor recoupling matrix is real orthogonal to numerical tolerance $10^{-10}$;
\item the anchor angular speed is nonzero.
\end{enumerate}
\end{definition}

The implementation contains $283$ ordered carriers in $24$ unordered external-spin families.  The independent-validation subset is constructed by first forming the label-$1$-to-$4$ family set and then retaining every carrier with a label $5$ or $6$ whose sorted external-label family is absent from that set.  This gives $208$ ordered carriers in $15$ independent families.

\section{Trigonometric quantum numbers and finite Racah data}
Let
\begin{equation}
q=e^{i\theta},
\qquad
\qnum{n}=\frac{\sin(n\theta)}{\sin\theta},
\qquad
\qfac{n}=\prod_{k=1}^{n}\qnum{k},
\qquad
\qfac{0}=1.
\label{eq:qnumber}
\end{equation}
The anchor is
\begin{equation}
\theta_0=\frac{\pi}{12}.
\label{eq:anchor}
\end{equation}

For an admissible triple $(a,b,e)$ in doubled labels, write
\begin{align}
\alpha&=\frac{a+b-e}{2},&
\beta&=\frac{a-b+e}{2},&
\gamma&=\frac{-a+b+e}{2},&
\sigma&=\frac{a+b+e}{2}+1.
\end{align}
The trigonometric triangle factor is
\begin{equation}
\Delta_q(a,b,e)
=
\left(
\frac{\qfac{\alpha}\qfac{\beta}\qfac{\gamma}}
{\qfac{\sigma}}
\right)^{1/2}.
\label{eq:qdelta}
\end{equation}
The square root is fixed at the positive real value at $\theta_0$ and then analytically continued inside the certified regular disk.

For six admissible doubled labels, define
\begin{align}
A_1&=\frac{a+b+e}{2},& A_2&=\frac{a+d+f}{2},&
A_3&=\frac{c+b+f}{2},& A_4&=\frac{c+d+e}{2},\\
B_1&=\frac{a+b+c+d}{2},&
B_2&=\frac{a+c+e+f}{2},&
B_3&=\frac{b+d+e+f}{2}.
\end{align}
The finite quantum Racah expression used in the software is
\begin{align}
\sixj{a}{b}{e}{c}{d}{f}
&=
\Delta_q(a,b,e)\Delta_q(a,d,f)
\Delta_q(c,b,f)\Delta_q(c,d,e)
\nonumber\\
&\quad\times
\sum_{z=\max A_i}^{\min B_j}
(-1)^z
\frac{\qfac{z+1}}
{\prod_{i=1}^{4}\qfac{z-A_i}\prod_{j=1}^{3}\qfac{B_j-z}}.
\label{eq:qracah}
\end{align}
No asymptotic replacement of this finite sum is used.

\section{The two-channel recoupling matrix}
For $E(J)=\{e_1,e_2\}$ and $G(J)=\{f_1,f_2\}$ in increasing order, define the raw matrix
\begin{equation}
\widetilde F_{ij}(\theta)
=
(-1)^{(a+b+c+d)/2}
\sqrt{\qnum{e_i+1}\qnum{f_j+1}}
\sixj{a}{b}{e_i}{c}{d}{f_j}.
\label{eq:rawF}
\end{equation}
A constant diagonal sign gauge is selected at $\theta_0$ so that the gauged matrix
\begin{equation}
F(\theta)=G_J\widetilde F(\theta)
\label{eq:gaugedF}
\end{equation}
has positive orientation and a fixed sign convention for its $(1,1)$ entry.  The gauge is frozen as $\theta$ varies.

\begin{assumption}[Regular orthogonal chamber]
For each declared carrier, the analytic continuation of \eqref{eq:gaugedF} is taken only in a simply connected disk on which all q-factorials in denominators are nonzero, all radicands in \eqref{eq:qdelta} are nonzero, and the square-root branches selected at $\theta_0$ remain analytic.
\end{assumption}

At real $\theta$ in the declared chamber, the matrix is orthogonal:
\begin{equation}
F(\theta)F(\theta)^T=I_2.
\label{eq:orthogonality}
\end{equation}
Differentiation gives
\begin{equation}
K(\theta):=F'(\theta)F(\theta)^T,
\qquad
K(\theta)^T=-K(\theta).
\label{eq:generator}
\end{equation}
Hence
\begin{equation}
K(\theta)=
\begin{pmatrix}
0&-\omega(\theta)\\
\omega(\theta)&0
\end{pmatrix}.
\label{eq:omega}
\end{equation}
The scalar $\omega$ is called the local angular speed of the recoupling frame.
